% GENERATED FILE. Edit canonical lessons, then run npm run generate:books.
% canonical-source-hash: fnv1a64:a2e9f464cf020c10
% canonical-lessons: HI-C549-vikalp, HI-C549-duvidha, HI-C549-asmanjas, HI-C549-tarjih, HI-C549-hicakna

\chapter{Two Ways to Go}
\label{ch:hi-two-ways-to-go}
\hlchaptermodality{549}

\begin{chapteropening}
\textbf{By the end of this chapter:} \emph{I can weigh two options and hesitate between them.}

Complete the last lesson of chapter 549: I can weigh two options and hesitate between them.
\end{chapteropening}

\section[\texorpdfstring{\hi{विकल्प} (vikalp)}{विकल्प (vikalp)}]{\hi{विकल्प} (vikalp) — an option; an alternative}

\label{lesson:HI-C549-vikalp}



\begin{quote}
Before the new one: say the Hindi for laziness, then the Hindi for to squander.
\end{quote}

\subsection*{You'll want to know: \hi{विकल्प}}
\textbf{\hi{विकल्प}} — \emph{vikalp} — \textquotedblleft{}an option; an alternative\textquotedblright{}.

\begin{quote}
\hi{हमारे} \hi{पास} \hi{दो} \hi{विकल्प} \hi{हैं}: \hi{बस} \hi{या} \hi{ट्रेन।}
\end{quote}

\begin{quote}
\emph{Hamāre pās do vikalp haiṁ: bas yā ṭren.}
\end{quote}

\begin{quote}
We have two options: the bus or the train.
\end{quote}

Masculine.

\begin{cousinweb}
\textbf{\hi{विकल्प}} (\emph{vikalp}) is a tatsama, taken whole from Sanskrit \emph{vikalpa}, \textquotedblleft{}alternation, option\textquotedblright{} (Monier-Williams).
\end{cousinweb}

\begin{grammarlens}[title={What you've built}]
One more word for saying what happens if.
\end{grammarlens}

\subsection*{Your turn}
\begin{itemize}
\raggedright
  \item \emph{Say it:} \emph{vikalp}
  \item \emph{Say it:} \emph{vikalp}, once more
  \item \emph{Say it:} \emph{Hamāre pās do vikalp haiṁ: bas yā ṭren.} — We have two options: the bus or the train.
  \item \emph{Recall:} say the Hindi for laziness, then the Hindi for to squander, then say \emph{vikalp} again
\end{itemize}

\begin{tcolorbox}[breakable,colback=teal!4,colframe=teal!35!black,title={Before you move on}]
What does \hi{विकल्प} mean? (\textquotedblleft{}An option; an alternative\textquotedblright{}.) Say it once more.
\end{tcolorbox}
\section[\texorpdfstring{\hi{दुविधा} (duvidhā)}{दुविधा (duvidhā)}]{\hi{दुविधा} (duvidhā) — a dilemma; being in two minds}

\label{lesson:HI-C549-duvidha}



\begin{quote}
Before the new one: say the Hindi for to squander, then the Hindi for an option.
\end{quote}

\subsection*{You'll want to know: \hi{दुविधा}}
\textbf{\hi{दुविधा}} — \emph{duvidhā} — \textquotedblleft{}a dilemma; being in two minds\textquotedblright{}.

\begin{quote}
\hi{दो} \hi{अच्छे} \hi{घर} \hi{देखकर} \hi{अरुण} \hi{दुविधा} \hi{में} \hi{पड़} \hi{गया।}
\end{quote}

\begin{quote}
\emph{Do acche ghar dekhkar Aruṇ duvidhā meṁ paṛ gayā.}
\end{quote}

\begin{quote}
After seeing two good houses, Arun was in two minds.
\end{quote}

Feminine.

\begin{cousinweb}
\textbf{\hi{दुविधा}} (\emph{duvidhā}) is usually traced to Sanskrit \emph{dvividha}, \textquotedblleft{}twofold, of two kinds\textquotedblright{} (Monier-Williams); the sense \textquotedblleft{}being of two minds\textquotedblright{} is the Hindi one.
\end{cousinweb}

\begin{grammarlens}[title={What you've built}]
One more word for saying what happens if.
\end{grammarlens}

\subsection*{Your turn}
\begin{itemize}
\raggedright
  \item \emph{Say it:} \emph{duvidhā}
  \item \emph{Say it:} \emph{duvidhā}, once more
  \item \emph{Say it:} \emph{Do acche ghar dekhkar Aruṇ duvidhā meṁ paṛ gayā.} — After seeing two good houses, Arun was in two minds.
  \item \emph{Recall:} say the Hindi for to squander, then the Hindi for an option, then say \emph{duvidhā} again
\end{itemize}

\begin{tcolorbox}[breakable,colback=teal!4,colframe=teal!35!black,title={Before you move on}]
What does \hi{दुविधा} mean? (\textquotedblleft{}A dilemma; being in two minds\textquotedblright{}.) Say it once more.
\end{tcolorbox}
\section[\texorpdfstring{\hi{असमंजस} (asmañjas)}{असमंजस (asmañjas)}]{\hi{असमंजस} (asmañjas) — indecision; uncertainty}

\label{lesson:HI-C549-asmanjas}



\begin{quote}
Before the new one: say the Hindi for an option, then the Hindi for a dilemma.
\end{quote}

\subsection*{You'll want to know: \hi{असमंजस}}
\textbf{\hi{असमंजस}} — \emph{asmañjas} — \textquotedblleft{}indecision; uncertainty\textquotedblright{}.

\begin{quote}
\hi{नौकरी} \hi{बदलने} \hi{को} \hi{लेकर} \hi{वह} \hi{असमंजस} \hi{में} \hi{है।}
\end{quote}

\begin{quote}
\emph{Naukrī badalne ko lekar vah asmañjas meṁ hai.}
\end{quote}

\begin{quote}
He cannot decide whether to change jobs.
\end{quote}

Masculine.

\begin{cousinweb}
\textbf{\hi{असमंजस}} (\emph{asmañjas}) is a tatsama, taken whole from Sanskrit \emph{asamañjasa}, \textquotedblleft{}unfit, unbecoming\textquotedblright{} (Monier-Williams). Hindi uses it for not knowing which way to turn.
\end{cousinweb}

\begin{grammarlens}[title={What you've built}]
One more word for saying what happens if.
\end{grammarlens}

\subsection*{Your turn}
\begin{itemize}
\raggedright
  \item \emph{Say it:} \emph{asmañjas}
  \item \emph{Say it:} \emph{asmañjas}, once more
  \item \emph{Say it:} \emph{Naukrī badalne ko lekar vah asmañjas meṁ hai.} — He cannot decide whether to change jobs.
  \item \emph{Recall:} say the Hindi for an option, then the Hindi for a dilemma, then say \emph{asmañjas} again
\end{itemize}

\begin{tcolorbox}[breakable,colback=teal!4,colframe=teal!35!black,title={Before you move on}]
What does \hi{असमंजस} mean? (\textquotedblleft{}Indecision; uncertainty\textquotedblright{}.) Say it once more.
\end{tcolorbox}
\section[\texorpdfstring{\hi{तरजीह} (tarjīh)}{तरजीह (tarjīh)}]{\hi{तरजीह} (tarjīh) — preference; priority}

\label{lesson:HI-C549-tarjih}



\begin{quote}
Before the new one: say the Hindi for a dilemma, then the Hindi for indecision.
\end{quote}

\subsection*{You'll want to know: \hi{तरजीह}}
\textbf{\hi{तरजीह}} — \emph{tarjīh} — \textquotedblleft{}preference; priority\textquotedblright{}.

\begin{quote}
\hi{दुकानदार} \hi{ने} \hi{पुराने} \hi{ग्राहकों} \hi{को} \hi{तरजीह} \hi{दी।}
\end{quote}

\begin{quote}
\emph{Dukāndār ne purāne grāhkoṁ ko tarjīh dī.}
\end{quote}

\begin{quote}
The shopkeeper gave priority to his old customers.
\end{quote}

Feminine.

\begin{cousinweb}
\textbf{\hi{तरजीह}} (\emph{tarjīh}) is usually traced to Arabic \emph{tarjīḥ}, \textquotedblleft{}giving more weight, preferring\textquotedblright{}, through Persian.
\end{cousinweb}

\begin{grammarlens}[title={What you've built}]
One more word for saying what happens if.
\end{grammarlens}

\subsection*{Your turn}
\begin{itemize}
\raggedright
  \item \emph{Say it:} \emph{tarjīh}
  \item \emph{Say it:} \emph{tarjīh}, once more
  \item \emph{Say it:} \emph{Dukāndār ne purāne grāhkoṁ ko tarjīh dī.} — The shopkeeper gave priority to his old customers.
  \item \emph{Recall:} say the Hindi for a dilemma, then the Hindi for indecision, then say \emph{tarjīh} again
\end{itemize}

\begin{tcolorbox}[breakable,colback=teal!4,colframe=teal!35!black,title={Before you move on}]
What does \hi{तरजीह} mean? (\textquotedblleft{}Preference; priority\textquotedblright{}.) Say it once more.
\end{tcolorbox}
\section[\texorpdfstring{\hi{हिचकना} (hicaknā)}{हिचकना (hicaknā)}]{\hi{हिचकना} (hicaknā) — to hesitate, to hold back}

\label{lesson:HI-C549-hicakna}



\begin{quote}
Before the new one: say the Hindi for indecision, then the Hindi for preference.
\end{quote}

\subsection*{You'll want to know: \hi{हिचकना}}
\textbf{\hi{हिचकना}} — \emph{hicaknā} — \textquotedblleft{}to hesitate, to hold back\textquotedblright{}.

\begin{quote}
\hi{बोलने} \hi{से} \hi{पहले} \hi{वह} \hi{थोड़ा} \hi{हिचका।}
\end{quote}

\begin{quote}
\emph{Bolne se pahle vah thoṛā hickā.}
\end{quote}

\begin{quote}
He hesitated a little before speaking.
\end{quote}

Past (he …): \hi{वह} \hi{हिचका}, \emph{vah hickā}.

\begin{cousinweb}
\textbf{\hi{हिचकना}} (\emph{hicaknā}) is a tadbhava, inherited from \emph{*hicc}, \textquotedblleft{}draw back\textquotedblright{}, a form CDIAL reconstructs from its descendants (CDIAL 14083).
\end{cousinweb}

\begin{grammarlens}[title={What you've built}]
One more word for saying what happens if.
\end{grammarlens}

\subsection*{Your turn}
\begin{itemize}
\raggedright
  \item \emph{Say it:} \emph{hicaknā}
  \item \emph{Say it:} \emph{hicaknā}, once more
  \item \emph{Say it:} \emph{Bolne se pahle vah thoṛā hickā.} — He hesitated a little before speaking.
  \item \emph{Recall:} say the Hindi for indecision, then the Hindi for preference, then say \emph{hicaknā} again
\end{itemize}

\begin{tcolorbox}[breakable,colback=teal!4,colframe=teal!35!black,title={Before you move on}]
What does \hi{हिचकना} mean? (\textquotedblleft{}To hesitate, to hold back\textquotedblright{}.) Say it once more.
\end{tcolorbox}
